\documentclass[aspectratio=169, xcolor=dvipsnames]{beamer}

% --- Theme Setup ---
\usetheme{Madrid}
\usecolortheme{default}

% --- Packages ---
\usepackage{graphicx}
\usepackage{xcolor}
\usepackage{booktabs}
\usepackage{hyperref}
\usepackage{adjustbox}
\usepackage{amssymb}
\usepackage{amsmath}
\usepackage{listings}
\usepackage{caption}
\usepackage{tikz}

\lstset{
  basicstyle=\ttfamily\small,
  keywordstyle=\color{blue}\bfseries,
  stringstyle=\color{red},
  showstringspaces=false,
  breaklines=true
}

% --- Title Information ---
\title{Module 49}
\subtitle{Concurrency Control/1}
\author{Partha Pratim Das}
\institute{Department of Computer Science and Engineering\\Indian Institute of Technology, Kharagpur\\ppd@cse.iitkgp.ac.in}
\date{\today}

\begin{document}

% Title Slide
\begin{frame}
    \titlepage
\end{frame}

\begin{frame}{Module Recap}
    \begin{itemize}[<+->]
        \item With proper planning, a database can be recovered back to a consistent state from inconsistent state in the face of system failures. Such a recovery is done via cascaded or cascadeless rollback
        \item View Serializability is a weaker serializability system for better concurrency. However, testing for view serializability is NP complete
    \end{itemize}
\end{frame}

\begin{frame}{Module Objectives}
    \begin{itemize}[<+->]
        \item Concurrency Control through design of serializable schedule is difficult in general. Hence we take a look into locking mechanism and Lock-Based Protocols
        \item We need to understand how locks may be implemented
    \end{itemize}
\end{frame}

\begin{frame}{Module Outline}
    \begin{itemize}[<+->]
        \item Concurrency Control
        \item Lock-Based Protocols
        \item Implementing Locking
    \end{itemize}
\end{frame}

\section{Concurrency Control}

\begin{frame}
  \vfill
  \centering
  \begin{beamercolorbox}[sep=8pt,center,shadow=true,rounded=true]{title}
    \usebeamerfont{title}Concurrency Control\par%
  \end{beamercolorbox}
  \vfill
\end{frame}

\begin{frame}{Concurrency Control}
    \begin{itemize}[<+->]
        \item A database must provide a mechanism that will ensure that all possible schedules are both:
        \begin{itemize}
            \item[$\triangleright$] Conflict serializable
            \item[$\triangleright$] Recoverable and, preferably, Cascadeless
        \end{itemize}
        \item A policy in which only one transaction can execute at a time generates serial schedules, but provides a poor degree of concurrency
        \item Concurrency-control schemes tradeoff between the amount of concurrency they allow and the amount of overhead that they incur
        \item Testing a schedule for serializability after it has executed is a little too late!
        \item Tests for serializability help us understand why a concurrency control protocol is correct
        \item \textbf{Goal}: To develop concurrency control protocols that will assure serializability
    \end{itemize}
\end{frame}

\begin{frame}{Concurrency Control (2)}
    \begin{itemize}[<+->]
        \item One way to ensure isolation is to require that data items be accessed in a \textbf{mutually exclusive manner}; that is, while one transaction is accessing a data item, no other transaction can modify that data item
        \item Should a transaction hold a lock on the whole database?
        \begin{itemize}
            \item[$\triangleright$] Would lead to strictly serial schedules - very poor performance
        \end{itemize}
        \item The most common method used to implement locking requirement is to allow a transaction to access a data item only if it is currently holding a lock on that item
    \end{itemize}
\end{frame}

\section{Lock-Based Protocols}

\begin{frame}
  \vfill
  \centering
  \begin{beamercolorbox}[sep=8pt,center,shadow=true,rounded=true]{title}
    \usebeamerfont{title}Lock-Based Protocols\par%
  \end{beamercolorbox}
  \vfill
\end{frame}

\begin{frame}{Lock-Based Protocols}
    \begin{block}{Lock}
        A lock is a mechanism to control concurrent access to a data item.
    \end{block}
    \begin{itemize}[<+->]
        \item Data items can be locked in two modes:
        \begin{itemize}
            \item[$\triangleright$] \textbf{exclusive (X) mode}: Data item can be both read as well as written. X-lock is requested using \texttt{lock-X} instruction
            \item[$\triangleright$] \textbf{shared (S) mode}: Data item can only be read. S-lock is requested using \texttt{lock-S} instruction
        \end{itemize}
        \item A transaction can unlock a data item Q by the \texttt{unlock(Q)} Instruction
        \item Lock requests are made to the concurrency-control manager by the programmer
        \begin{itemize}
            \item[$\triangleright$] Transaction can proceed only after request is granted
        \end{itemize}
    \end{itemize}
\end{frame}

\begin{frame}{Lock-Based Protocols (2): Lock Compatibility Matrix}
    \begin{block}{Lock-Compatibility Matrix}
        A lock compatibility matrix is used which states whether a data item can be locked by two transactions at the same time.
    \end{block}
    \vspace{1em}
    \uncover<2->{
        \begin{center}
            \begin{tabular}{|l|l|l|}
            \hline
            & \textbf{Shared} & \textbf{Exclusive} \\ \hline
            \textbf{Shared} & Yes & No \\ \hline
            \textbf{Exclusive} & No & No \\ \hline
            \end{tabular}
        \end{center}
    }
\end{frame}

\begin{frame}{Lock-Based Protocols (3)}
    \begin{itemize}[<+->]
        \item \textbf{Requesting for / Granting of a Lock}
        \begin{itemize}
            \item[$\triangleright$] A transaction may be granted a lock on an item if the requested lock is compatible with locks already held on the item by other transactions
        \end{itemize}
        \item \textbf{Sharing a Lock}
        \begin{itemize}
            \item[$\triangleright$] Any number of transactions can hold shared locks on an item
            \item[$\triangleright$] But if any transaction holds an exclusive lock on the item no other transaction may hold any lock on the item
        \end{itemize}
        \item \textbf{Waiting for a Lock}
        \begin{itemize}
            \item[$\triangleright$] If a lock cannot be granted, the requesting transaction is made to wait till all incompatible locks held by other transactions have been released
        \end{itemize}
    \end{itemize}
\end{frame}

\begin{frame}{Lock-Based Protocols (4)}
    \begin{itemize}[<+->]
        \item \textbf{Holding a Lock}
        \begin{itemize}
            \item[$\triangleright$] A transaction must hold a lock on a data item as long as it accesses that item
        \end{itemize}
        \item \textbf{Unlocking / Releasing a Lock}
        \begin{itemize}
            \item[$\triangleright$] Transaction $T_i$ may unlock a data item that it had locked at some earlier point
            \item[$\triangleright$] It is not necessarily desirable for a transaction to unlock a data item immediately after its final access of that data item, since serializability may not be ensured
        \end{itemize}
    \end{itemize}
\end{frame}

\begin{frame}[fragile]{Lock-Based Protocols: Example: Serial Schedule}
    \begin{itemize}[<+->]
        \item Let A and B be two accounts that are accessed by transactions $T_1$ and $T_2$.
        \item Transaction $T_1$ transfers \$50 from account B to account A
        \item Transaction $T_2$ displays the total amount of money in accounts A and B
        \item Suppose that the values of accounts A and B are \$100 and \$200, respectively
        \item If executed serially (e.g. $T_1$, $T_2$), $T_2$ will display \$300
    \end{itemize}
        \begin{columns}[T]
            \begin{column}<6->{0.48\textwidth}
                \textbf{T1:}
                \begin{lstlisting}
lock-X(B);
read(B); 
B := B - 50; 
write(B); 
unlock(B); 
lock-X(A); 
read(A); 
A := A + 50; 
write(A); 
unlock(A);
                \end{lstlisting}
            \end{column}
            \begin{column}{0.48\textwidth}
                \textbf{T2:}
                \begin{lstlisting}
lock-S(A); 
read(A); 
unlock(A); 
lock-S(B); 
read(B); 
unlock(B); 
display(A + B);
                \end{lstlisting}
            \end{column}
        \end{columns}
\end{frame}

\begin{frame}[fragile]{Lock-Based Protocols: Example: Concurrent Schedule: Bad}
    \begin{alertblock}{Problem: Unlocking Too Early}
        In this case, transaction $T_2$ displays \$250, which is incorrect. The reason for this mistake is that the transaction $T_1$ unlocked data item B too early, as a result of which $T_2$ saw an inconsistent state.
    \end{alertblock}
        \begin{columns}[T]
            \begin{column}<2->{0.48\textwidth}
                \textbf{T1:}
                \begin{lstlisting}
lock-X(B);
read(B); 
B := B - 50; 
write(B); 
unlock(B);


lock-X(A); 
read(A); 
A := A + 50; 
write(A); 
unlock(A);
                \end{lstlisting}
            \end{column}
            \begin{column}{0.48\textwidth}
                \textbf{T2:}
                \begin{lstlisting}





lock-S(A); 
read(A); 
unlock(A); 
lock-S(B); 
read(B); 
unlock(B); 
display(A + B);
                \end{lstlisting}
            \end{column}
        \end{columns}
    \uncover<3->{
        \textbf{Solution:} Delay unlocking till the end
    }
\end{frame}

\begin{frame}{Lock-Based Protocols: Example: Concurrent Schedule: Deadlock}
    \begin{itemize}[<+->]
        \item Since $T_3$ is holding an exclusive mode lock on B and $T_4$ is requesting a shared-mode lock on B, $T_4$ is waiting for $T_3$ to unlock B
        \item Similarly, since $T_4$ is holding a shared-mode lock on A and $T_3$ is requesting an exclusive-mode lock on A, $T_3$ is waiting for $T_4$ to unlock A
    \end{itemize}
    \uncover<3->{
        \begin{alertblock}{Deadlock}
            We have arrived at a state where neither of these transactions can ever proceed with its normal execution. This situation is called deadlock.
        \end{alertblock}
    }
    \begin{itemize}[<+->]
        \item When deadlock occurs, the system must roll back one of the two transactions.
        \item Once a transaction has been rolled back, the data items that were locked by that transaction are unlocked and available to the other transaction.
    \end{itemize}
\end{frame}

\begin{frame}[fragile]{Deadlocks (Cont'd)}
    \begin{exampleblock}{Example of Deadlock}
        \begin{columns}[T]
            \begin{column}{0.48\textwidth}
                \textbf{T3:}
                \begin{lstlisting}
lock-X(B);
read(B); 
B := B - 50; 
write(B); 

lock-X(A); 
                \end{lstlisting}
            \end{column}
            \begin{column}{0.48\textwidth}
                \textbf{T4:}
                \begin{lstlisting}

lock-S(A); 

read(A); 
lock-S(B);
                \end{lstlisting}
            \end{column}
        \end{columns}
    \end{exampleblock}
\end{frame}

\begin{frame}{Starvation}
    \begin{itemize}[<+->]
        \item In addition to deadlocks, there is a possibility of \textbf{Starvation}
        \item Starvation occurs if the concurrency control manager is badly designed. For example:
        \begin{itemize}
            \item[$\triangleright$] A transaction may be waiting for an X-lock on an item, while a sequence of other transactions request and are granted an S-lock on the same item
            \item[$\triangleright$] The same transaction is repeatedly rolled back due to deadlocks
        \end{itemize}
        \item Concurrency control manager can be designed to prevent starvation
        \item Starvation is also loosely referred to as \textbf{Livelock}
    \end{itemize}
\end{frame}

\begin{frame}{Cascading Rollback}
    \begin{itemize}[<+->]
        \item The potential for deadlock exists in most locking protocols. Deadlocks are a necessary evil.
        \item When a deadlock occurs there is a possibility of \textbf{cascading roll-backs}
        \item Cascading roll-back is possible under two-phase locking
        \item In the schedule below, the failure of $T_5$ after the \texttt{read(A)} step of $T_7$ leads to cascading rollback of $T_6$ and $T_7$.
    \end{itemize}
    \uncover<5->{
        \begin{center}
            \begin{adjustbox}{max width=\textwidth}
            \begin{tabular}{|l|l|l|}
            \hline
            \textbf{$T_5$} & \textbf{$T_6$} & \textbf{$T_7$} \\ \hline
            \texttt{lock-X(A)} & & \\ 
            \texttt{read(A)} & & \\ 
            \texttt{lock-S(B)} & & \\ 
            \texttt{read(B)} & & \\ 
            \texttt{write(A)} & & \\ 
            \texttt{unlock(A)} & \texttt{lock-X(A)} & \\ 
            & \texttt{read(A)} & \\ 
            & \texttt{write(A)} & \\ 
            & \texttt{unlock(A)} & \texttt{lock-S(A)} \\ 
            & & \texttt{read(A)} \\ \hline
            \end{tabular}
            \end{adjustbox}
        \end{center}
    }
\end{frame}

\begin{frame}{Two-Phase Locking Protocol}
    \begin{block}{Two-Phase Locking}
        This protocol ensures conflict-serializable schedules by dividing execution into two phases.
    \end{block}
    \begin{itemize}[<+->]
        \item \textbf{Phase 1: Growing Phase}
        \begin{itemize}
            \item[$\triangleright$] Transaction may obtain locks
            \item[$\triangleright$] Transaction may not release locks
        \end{itemize}
        \item \textbf{Phase 2: Shrinking Phase}
        \begin{itemize}
            \item[$\triangleright$] Transaction may release locks
            \item[$\triangleright$] Transaction may not obtain locks
        \end{itemize}
        \item The protocol assures serializability. It can be proved that the transactions can be serialized in the order of their lock points (i.e., the point where a transaction acquired its final lock).
    \end{itemize}
\end{frame}

\begin{frame}{Two-Phase Locking Protocol (2)}
    \begin{itemize}[<+->]
        \item There can be conflict serializable schedules that cannot be obtained if two-phase locking is used
        \item However, in the absence of extra information (that is, ordering of access to data), two-phase locking is needed for conflict serializability in the following sense:
        \begin{itemize}
            \item[$\triangleright$] Given a transaction $T_i$ that does not follow two-phase locking, we can find a transaction $T_j$ that uses two-phase locking, and a schedule for $T_i$ and $T_j$ that is not conflict serializable
        \end{itemize}
    \end{itemize}
\end{frame}

\begin{frame}{Lock Conversions}
    \begin{itemize}[<+->]
        \item \textbf{Two-phase locking with lock conversions:}
        \item \textbf{First Phase:}
        \begin{itemize}
            \item[$\triangleright$] can acquire a \texttt{lock-S} on item
            \item[$\triangleright$] can acquire a \texttt{lock-X} on item
            \item[$\triangleright$] can convert a \texttt{lock-S} to a \texttt{lock-X} (upgrade)
        \end{itemize}
        \item \textbf{Second Phase:}
        \begin{itemize}
            \item[$\triangleright$] can release a \texttt{lock-S}
            \item[$\triangleright$] can release a \texttt{lock-X}
            \item[$\triangleright$] can convert a \texttt{lock-X} to a \texttt{lock-S} (downgrade)
        \end{itemize}
        \item This protocol assures serializability. But still relies on the programmer to insert the various locking instructions.
    \end{itemize}
\end{frame}

\begin{frame}[fragile]{Automatic Acquisition of Locks: Read}
    \begin{itemize}[<+->]
        \item A transaction $T_i$ issues the standard read/write instruction, without explicit locking calls
        \item The operation \texttt{read(D)} is processed as:
    \end{itemize}
    \begin{onlyenv}<3->
        \begin{lstlisting}[language=SQL]
if Ti has a lock on D
   then read(D)
else begin 
   if necessary, wait until no other transaction has a lock-X on D
   grant Ti a lock-S on D; 
   read(D)
end
        \end{lstlisting}
    \end{onlyenv}
\end{frame}

\begin{frame}[fragile]{Automatic Acquisition of Locks: Write}
    \begin{itemize}[<+->]
        \item The operation \texttt{write(D)} is processed as:
    \end{itemize}
    \begin{onlyenv}<2->
        \begin{lstlisting}[language=SQL]
if Ti has a lock-X on D
   then write(D)
else begin 
   if necessary, wait until no other transaction has any lock on D, 
   if Ti has a lock-S on D
       then upgrade lock on D to lock-X
       else grant Ti a lock-X on D
   write(D)
end;
        \end{lstlisting}
    \end{onlyenv}
    \begin{itemize}[<+->]
        \item<3-> All locks are released after commit or abort
    \end{itemize}
\end{frame}

\begin{frame}{More Two Phase Locking Protocols}
    \begin{itemize}[<+->]
        \item To avoid Cascading roll-back, follow a modified protocol called \textbf{strict two-phase locking}
        \begin{itemize}
            \item[$\triangleright$] a transaction must hold all its exclusive locks till it commits/aborts
        \end{itemize}
        \item \textbf{Rigorous two-phase locking} is even stricter
        \begin{itemize}
            \item[$\triangleright$] All locks are held till commit/abort. In this protocol transactions can be serialized in the order in which they commit
        \end{itemize}
        \item Note that concurrency goes down as we move to more and more strict locking protocol
    \end{itemize}
\end{frame}

\section{Implementation of Locking}

\begin{frame}
  \vfill
  \centering
  \begin{beamercolorbox}[sep=8pt,center,shadow=true,rounded=true]{title}
    \usebeamerfont{title}Implementation of Locking\par%
  \end{beamercolorbox}
  \vfill
\end{frame}

\begin{frame}{Implementation of Locking}
    \begin{itemize}[<+->]
        \item A lock manager can be implemented as a separate process to which transactions send lock and unlock requests
        \item The lock manager replies to a lock request by sending a lock grant messages (or a message asking the transaction to roll back, in case of a deadlock)
        \item The requesting transaction waits until its request is answered
        \item The lock manager maintains a data-structure called a \textbf{lock table} to record granted locks and pending requests
        \item The lock table is usually implemented as an in-memory hash table indexed on the name of the data item being locked
    \end{itemize}
\end{frame}

\begin{frame}{Lock Table}
    \begin{itemize}[<+->]
        \item Lock table also records the type of lock granted or requested
        \item New request is added to the end of the queue of requests for the data item, and granted if it is compatible with all earlier locks
        \item Unlock requests result in the request being deleted, and later requests are checked to see if they can now be granted
        \item If transaction aborts, all waiting or granted requests of the transaction are deleted
        \item Lock manager may keep a list of locks held by each transaction, to implement this efficiently
    \end{itemize}
\end{frame}

\begin{frame}{Module Summary}
    \begin{itemize}[<+->]
        \item Understood the locking mechanism and protocols
        \item Realized that deadlock is a peril of locking and needs to be handled through rollback
    \end{itemize}
    \vspace{2em}
    \uncover<3->{
        \begin{center}
        \small \textit{Slides used in this presentation are borrowed from \url{http://db-book.com/} with kind permission of the authors.} \\
        \small \textit{Edited and new slides are marked with 'PPD'.}
        \end{center}
    }
\end{frame}

\end{document}
